\documentclass[17pt]{flash}
\startdatakeys
\cycle{3}
\connaissancesRequises{}
\competencesTravaillees{}
\enddatakeys
\begin{document}

\flashvisibletimertrue
\FlashSetItemSep{40pt}
\FlashSetVtop{-2em}
\FlashSetTitle{Volume d'un pavé droit}

\FlashShowConsignes{%
    \item Calculer le volume de chaque pavé droit.
    \item Écrire le résultat avec l'unité adaptée.
}

\flashlist{
    \flashitem{20}{Un pavé droit mesure $3$ cm de long, $4$ cm de large et $5$ cm de haut.}
    \flashitem{20}{Un cube a des arêtes de $2$ cm.}
    \flashitem{30}{Un pavé droit mesure $10$ cm de long, $2$ cm de large et $3$ cm de haut.}
    \flashitem{30}{La base d'un pavé droit a une aire de $6$ cm$^2$ et sa hauteur est de $4$ cm.}
}

\end{document}
